\subsection{RINGS}

DEFINITION 1. A ring is a set A with two laws of composition called respectively addition and multiplication, satisfying the following axioms:

(AN I) Under addition A is a commutative group.

(AN II) Multiplication is associative and possesses an identity element.

(AN III) Multiplication is distributive with respect to addition.

The ring A is said to be commutative if its multiplication is commutative.

In what follows, \((x,y)\mapsto x+y\) denotes addition and \((x,y)\mapsto xy\) multiplication; 0 denotes the identity element for addition and 1 that for multiplication. Finally, -x denotes the negative of x under addition. The axioms of a ring are therefore expressed by the following identities:

\[
x + (y + z) = (x + y) + z \quad (\text { associativity   of   addition })\tag{1}
\]

\[
x + y = y + x \quad (\text { commutativity   of   addition })\tag{2}
\]

\[
0 + x = x + 0 = x \quad (\text { zero })\tag{3}
\]

\[
x + (- x) = (- x) + x = 0 \quad (\text { negative })\tag{4}
\]

\[
x (y z) = (x y) z \quad (\text { associativity   of   multiplication })\tag{5}
\]

\[
x. 1 = 1. x = x \quad (\text { unit   element })\tag{6}
\]

\[
\left. \begin{array}{l} (x + y). z = x z + y z \\ x. (y + z) = x y + x z \end{array} \right\} (\text {distributivity})\tag{7, 8}
\]

Finally the ring A is commutative if \(xy = yx\) for \(x, y\) in A.

With addition alone A is a commutative group called the additive group of A. For all \(x \in A\) , we define the left homothety \(\gamma_{x}\) and right homothety \(\delta_{x}\) by \(\gamma_{x}(y) = xy\) , \(\delta_{x}(y) = yx\) . By formulae (7) and (8), \(\gamma_{x}\) and \(\delta_{x}\) are endomorphisms of the additive group of A and thus map zero to zero and negative to negative. Therefore

\[
x. 0 = 0. x = 0\tag{9}
\]

\[
x. (- y) = (- x). y = - x y;\tag{10}
\]

it follows that \((-x)(-y) = -((-x).y) = -(-xy)\) , whence

\[
(- x) (- y) = x y.\tag{11}
\]

Formulae (10) and (11) constitute the sign rule. It follows that

\[
- x = (- 1) x = x (- 1)
\]

and \((-1)(-1) = 1\) .

From (11) it follows by induction on \(n\) that

\[
(- x) ^ {n} = \left\{ \begin{array}{l l} x ^ {n} & \text { if   } n \text {   is   even } \\ - x ^ {n} & \text { if   } n \text {   is   odd. } \end{array} \right.\tag{12}
\]

When we speak of cancellable elements, invertible elements, permutable elements, central elements, centralizer or centre of a ring A, all these notions will refer to the multiplication on A. If \(x, y \in A\) and y is invertible, the element \(xy^{-1}\) of A is also denoted by x/y when A is commutative. The set of invertible elements of A is stable under multiplication. Under the law induced by multiplication it is a group called the multiplicative group of A, sometimes denoted by A*.

Let \(x, y\) be in A. \(x\) is said to be a left (resp. right) multiple of \(y\) if there exists \(y' \in \mathbf{A}\) such that \(x = y'y\) (resp. \(x = yy'\) ); it is also said that \(y\) is a right (resp.

left) divisor of x. When A is commutative, there is no need to distinguish between ``left'' and ``right''.

In conformity with the above terminology, every element \(y \in A\) would be considered as a right and left divisor of 0; but, by an abuse of language, in general the term ``right (resp. left) divisor of 0'' is reserved for elements y such that there exists \(x \neq 0\) in A satisfying the relation xy = 0 (resp. yx = 0). In other words, the right (resp. left) divisors of zero are the right (resp. left) non-cancellable elements.

Let \(x \in A\) . \(x\) is called nilpotent if there exists an integer \(n > 0\) with \(x^n = 0\) . The element \(1 - x\) is then invertible, with inverse equal to

\[
1 + x + x ^ {2} + \dots + x ^ {n - 1}.
\]

As A is a commutative group under addition, the element \(nx\) for \(n \in \mathbf{Z}\) and \(x \in \mathbf{A}\) has been defined (§ 2, no. 8). As \(\gamma_x\) and \(\delta_x\) are endomorphisms of the additive group \(\mathbf{A}\) , \(\gamma_x(ny) = n\gamma_x(y)\) and \(\delta_y(nx) = n\delta_y(x)\) , whence

\[
x. (n y) = (n x). y = n. (x y).
\]

In particular, \(nx = (n.1)x\) .

A set A with addition and multiplication satisfying the axioms of a ring with the exception of that assuring the existence of the identity element under multiplication, is called a pseudo-ring.

